% !TEX root = ../algebra_02.tex

\section{Задание группы образующими и определяющими соотношениями}

\begin{Definition}{Нормальное замыкание}{}
	Пусть $G$ — группа, $M \subset G$.
	Нормальным замыканием множества $M$ (или нормальной подгруппой, порожденной $M$) называется:
	\[ \langle\langle M \rangle\rangle := \bigcap_{\substack{M \subset H \\ H \triangleleft G}} H \]
	Очевидно, что по построению $\langle\langle M \rangle\rangle \triangleleft G$.
\end{Definition}

\begin{Definition}{Задание группы образующими и соотношениями}{}
	Запись вида:
	\[ \langle f_1, \dots, f_n \mid r_1, \dots, r_m \rangle \]
	где $r_1, \dots, r_m \in F_n$, обозначает факторгруппу:
	\[ F_n / \langle\langle r_1, \dots, r_m \rangle\rangle \]
	Данная конструкция называется группой, заданной образующими $f_1, \dots, f_n$ и определяющими соотношениями $r_1, \dots, r_m$.
\end{Definition}

\begin{Remark}{Табличное копредставление конечной группы}{}
	Любую конечную группу можно задать образующими и соотношениями.
	Пусть $|G| = n$, а её элементы занумерованы как $G = \{g_1, \dots, g_n\}$. Очевидно, что $G = \langle g_1, \dots, g_n \rangle$.
	Рассмотрим свободную группу $F_n$ со свободными образующими $f_1, \dots, f_n$. Каждому равенству $g_i g_j = g_k$, выполняющемуся в группе $G$, сопоставим соотношение $f_i f_j f_k^{-1}$. Тогда группу $G$ можно представить как:
	\[ G \cong \langle f_1, \dots, f_n \mid \{ f_i f_j f_k^{-1} \}_{i,j,k=1}^n \text{ для всех } i,j,k \text{ таких, что } g_i g_j = g_k \text{ в } G \rangle \]
\end{Remark}

\begin{Example}{Копредставление группы диэдра $D_5$}{}
	Покажем, что $D_5 \cong \langle S, T \mid S^2, T^5, (ST)^2 \rangle$.
	Пусть $F_2$ — свободная группа со свободными образующими $S, T$. Зададим гомоморфизм $\varphi: F_2 \to D_5$, действующий на образующие следующим образом:
	\[ S \mapsto S_0 \text{ (симметрия)}, \quad T \mapsto R_1 \text{ (поворот)} \]
	В группе $D_5$ элементы $S_0$ и $R_1$ удовлетворяют соотношениям $S_0^2 = e$, $R_1^5 = e$ и $(S_0 R_1)^2 = e$.
	Следовательно:
	\[ S^2 \in \ker \varphi, \quad T^5 \in \ker \varphi, \quad (ST)^2 \in \ker \varphi \]
	Определим нормальную подгруппу $N := \langle\langle S^2, T^5, STST \rangle\rangle$.
	Из включений выше очевидно, что $N \subset \ker \varphi$.

	Рассмотрим факторгруппу $F_2 / N = \langle s, t \rangle$, где $s = SN$, $t = TN$. По определению факторгруппы по нормальному замыканию, в $F_2/N$ справедливы тождества: $ s^2 = e, \quad t^5 = e, \quad stst = e $
	Из последнего равенства выразим $st$: $st = t^{-1}s^{-1} = t^4 s $
	Используя это соотношение (позволяющее перебрасывать $s$ вправо от $t$), любой элемент $g \in F_2 / N$ можно привести к виду $t^k s^l$:
	\[ g = s^{i_1} t^{j_1} s^{i_2} t^{j_2} \dots = t^k s^l \]
	где показатели степеней ограничены: $0 \le k \le 4$ и $0 \le l \le 1$. Следовательно, порядок факторгруппы ограничен сверху: $|F_2 / N| \le 10$.

	По теореме о гомоморфизме, так как $N \subset \ker \varphi$, существует индуцированный гомоморфизм $\widetilde{\varphi}: F_2 / N \to D_5$.
	Поскольку $\langle S_0, R_1 \rangle = D_5$, исходный гомоморфизм $\varphi$ является сюръективным, а значит, и $\widetilde{\varphi}$ сюръективен. Мы имеем сюръективное отображение группы, в которой не более 10 элементов, на группу из ровно 10 элементов ($|D_5| = 10$). Это возможно только в том случае, если $|F_2 / N| = 10$ и $\widetilde{\varphi}$ является биекцией. Таким образом, $\widetilde{\varphi}$ — изоморфизм.
\end{Example}

\begin{Example}{Копредставление группы $(\mathbb{Z}/6\mathbb{Z}) \times (\mathbb{Z}/10\mathbb{Z})$}{}
	Зададим образующими и соотношениями группу $G = (\mathbb{Z}/6\mathbb{Z}) \times (\mathbb{Z}/10\mathbb{Z})$.

	Эта группа порождается двумя элементами: $a = (\bar{1}, \bar{0})$ и $b = (\bar{0}, \bar{1})$. Порядки этих элементов равны $6$ и $10$ соответственно, то есть $a^6 = e$ и $b^{10} = e$. Так как группа является прямым произведением абелевых групп, она сама абелева, следовательно её образующие коммутируют: $ab = ba$, что эквивалентно соотношению $aba^{-1}b^{-1} = e$.

	Рассмотрим факторгруппу $\widetilde{G} = \langle a, b \mid a^6, b^{10}, aba^{-1}b^{-1} \rangle$. Любой элемент этой группы благодаря коммутативности образующих представим в виде $a^i b^j$. Из соотношений $a^6=e$ и $b^{10}=e$ следует, что достаточно рассматривать показатели $0 \le i < 6$ и $0 \le j < 10$. Таким образом, порядок факторгруппы $|\widetilde{G}| \le 6 \times 10 = 60$.

	С другой стороны, существует очевидный сюръективный гомоморфизм из свободной группы $F_2$ в $G$, при котором образующие переходят в $a$ и $b$. Ядро этого гомоморфизма содержит нормальное замыкание $\langle\langle a^6, b^{10}, aba^{-1}b^{-1} \rangle\rangle$. По теореме о гомоморфизме существует сюръекция $\widetilde{G} \to G$. Поскольку $|G| = 60$, а $|\widetilde{G}| \le 60$, отображение является изоморфизмом.

	Итоговое копредставление:
	\[ (\mathbb{Z}/6\mathbb{Z}) \times (\mathbb{Z}/10\mathbb{Z}) \cong \langle a, b \mid a^6, b^{10}, aba^{-1}b^{-1} \rangle \]
\end{Example}
